% 原文 pp.81–90（印刷 pp.1199–1208）
\section{Lecture 6：Bose–Fermi 对应及其应用}

\begin{keytakeaways}
本讲在二维（欧氏）时空中建立含费米子的 $U(1)$ 规范理论与自由玻色子之间的\textbf{玻色化}图像：无质量费米子的手征流代数与半径 $R=1/\sqrt2$ 的角度值玻色场的顶点算子代数等同，其定量表达是算子对应 $\xi(\psi_\pm)=e^{\pm i\phi_\pm}$ 与 Jacobi 三重积恒等式所给出的配分函数恒等式。所用数学工具是 Spin 丛上 Dirac 算子的指标、Pfaffian 线丛的曲率（它给出左右表示二次指标相等这一反常抵消条件）以及 Riemann 面上的整体路径积分恒等式。主要结论有三：手征对称性因测度变换带上因子 $e^{i\gamma I}$ 而反常，从而 $\theta$ 角可被吸收进复质量 $\tilde m=me^{-i\theta}$；四费米子耦合理论在玻色变量下是自由的（$\beta$ 函数为零），而 sine-Gordon 模型在 $R=1/\sqrt2$ 时等价于自由有质量费米子（massive Thirring 模型）；把该对应用于规范理论，可在 $m_*$ 小时得到显式的有效势 $U(\phi)$，其极小在 $\theta=\pi$ 处的行为与 Ising 型临界点相联系。
\end{keytakeaways}

% ===== 原文 p.81 =====
\subsection{含费米子的二维规范理论}
% 原文 §6.1

今天我们考虑含费米子的二维规范理论。我们将使用欧氏（Euclidean）拉格朗日量。

在 2 维时，两个旋量表示 $S_+,S_-$ 都是一维且复的。它们由公式 $\theta\mapsto e^{\pm i\theta/2}$，$\theta\in[0,4\pi)$ 定义，其中 $\theta$ 是 $\mathrm{Spin}(2)$ 上的参数。

设 $\Sigma$ 是一张给定了 Spin 结构（Spin structure）的 Riemann 面。于是 $S_+,S_-$ 定义了 $\Sigma$ 上的全纯线丛，它们被称为 Spin 丛（Spin bundles），并用同一个字母记之。我们有 $S_+\otimes S_-=\mathbb{C}$，以及 $S_+^2=T\Sigma$、$S_-^2=T^*\Sigma$。$\Pi S_-$（相应地 $\Pi S_+$）的截面将被称为左行（相应地右行）费米子，这是与 Minkowski 图像类比而得名的。

我们用 $D_+:S_-\to S_+$、$D_-:S_+\to S_-$ 记 Spin 丛中相应的 Dirac 算子。

我们将做规范群为 $G$（一个紧 Lie 群）的规范理论。设 $E_L,E_R$ 是 $G$ 的正交、幺模（unimodular）表示。设 $P$ 是 $\Sigma$ 上的主 $G$-丛，并用 $E_R,E_L$ 记与表示 $E_R,E_L$ 相伴的正交向量丛。设 $\psi_+,\psi_-$ 分别是丛 $E_R\otimes S_+$、$E_L\otimes S_-$ 的截面。

含费米子的二维规范理论的拉格朗日量为

\begin{equation}
  \tag{6.1}
  \label{eq:6-1}
  \mathcal{L}=\int d^2x\left(\frac{1}{4e^2}|\ast F_A|^2+\frac{1}{4\pi}\overrightarrow{\bar\psi}_+\,\overrightarrow{D^A_-}\,\overrightarrow{\psi}_++\frac{1}{4\pi}\overrightarrow{\bar\psi}_-\,\overrightarrow{D^A_+}\,\overrightarrow{\psi}_-\right),
\end{equation}

其中 $A$ 是 $P$ 中的联络，而 $D^A_\pm$ 是相应的 Dirac 算子。

% [Tutor Note: (6.1) 中费米子动能项的系数是 $1/4\pi$，而一旦按 (6.4) 把 $\mathbb{R}$ 上的实二维表示 $E$ 分解为复分量 $(\psi_\pm,\bar\psi_\pm)$，同一项的系数就变成 $1/2\pi$。这不是笔误：实表示的两个实分量各贡献一份动能。若在此处误把 $E$ 当作复一维表示，则 (6.4) 之后的玻色化系数（$|d\phi|^2$ 的 $(1+g)/8\pi$ 与质量项的 $m/2\pi$）会整体差一个因子 2。另注意箭头的含义：$\overrightarrow{\bar\psi}$ 是 (6.4) 下方所定义的“行向量”记号，见 p.82。]

为了使这个理论在量子力学上说得通，表示 $E_L,E_R$ 还必须满足一个附加条件。为导出这个条件，回顾 (6.1) 的配分函数由下式给出

\begin{equation}
  \tag{6.2}
  \label{eq:6-2}
  Z=\int DA\int D\overrightarrow{\bar\psi}_+\,D\overrightarrow{\bar\psi}_-\,e^{-\mathcal{L}}.
\end{equation}

（在此积分中，我们对所有拓扑类型的主丛求和）。费米子积分容易计算：它等于 $I_A=\mathrm{Pf}(D^A_-|_{E_R})\mathrm{Pf}(D^A_+|_{E_L})$，其中\footnote{原文页脚注：\emph{Notes by Pavel Etingof and David Kazhdan}——本讲笔记由 Pavel Etingof 与 David Kazhdan 整理。}

% ===== 原文 p.82 =====
$\mathrm{Pf}$ 表示 Pfaffian（普法夫式）。表达式 $I_A$ 是 Pfaffian 线丛 $B=\mathrm{Pf}(D^A_-|_{E_R})\mathrm{Pf}(D^A_+|_{E_L})$ 在联络的规范等价类空间上的一个截面。为使对 $A$ 的积分有意义，这个表达式应当是一个函数，亦即丛 $B$ 必须是平凡的。可以证明（利用 Pfaffian 线丛上联络曲率的标准数学公式），这归结为如下条件

\begin{equation}
  \tag{6.3}
  \label{eq:6-3}
  \operatorname{Tr}\big(\rho_L(t)\rho_L(t')\big)=\operatorname{Tr}\big(\rho_R(t)\rho_R(t')\big),
\end{equation}

其中 $t,t'\in\mathfrak{g}$，$\mathfrak{g}$ 是规范群 $G$ 的 Lie 代数，而 $\rho_{L,R}:\mathfrak{g}\to SO(E_{L,R})$ 是表示映射。若 $G$ 单，且 $E_L,E_R$ 不可约，则此条件意味着 $E_L,E_R$ 的 Casimir 相同。方程 (6.3) 称为反常抵消（anomaly cancellation）条件。从另一个观点看，它（在四维情形下）已在 Lecture 3.7 中描述过。

% [Tutor Note: 作者说“可以证明”而省略的，是 Pfaffian 线丛上 Quillen 型联络的曲率公式：其曲率正比于 $\operatorname{Tr}\rho^2(F)$ 的积分，因此 $B$ 的平凡性要求左右两个表示的二次指标相等，即 (6.3)。注意 (6.3) 只对 $\mathfrak{g}$ 中元素成立、并未要求 $E_L\cong E_R$ 作为表示同构；对 $G=U(1)$ 它退化为“左右荷的平方相等”，故 $E_L=E_R$ 只是一维情形下的结论。]

今天我们将考虑一个简单情形：$G=U(1)$，且 $E_L,E_R$ 是不可约的二维实表示。此时 (6.3) 说明 $E_L,E_R$ 相同：$E_L=E_R=E$。我们将取 $E=\mathbb{C}$，度量为 $|z|^2$，$U(1)$ 通过乘法作用（但要记住张量积 $S_\pm\otimes E$ 是在 $\mathbb{R}$ 上取的）。我们把 $S_\pm\otimes E$ 分解为两个二维表示的直和。利用这个分解，我们将写 $\overrightarrow{\bar\psi}_\pm=(\psi_\pm,\bar\psi_\pm)$，其中对 $z\in U(1)$ 有 $z(\psi,\bar\psi)=(z\psi,\bar z\bar\psi)$。

在新记号下，拉格朗日量 (6.1) 具有形式

\begin{equation}
  \tag{6.4}
  \label{eq:6-4}
  \mathcal{L}=\int d^2x\left(\frac{1}{4e^2}|\ast F_A|^2+\frac{1}{2\pi}\bar\psi_+D^A_-\psi_++\frac{1}{2\pi}\bar\psi_-D^A_+\psi_-\right),
\end{equation}

我们还可以给这个拉格朗日量加上一个拓扑项：

\begin{equation}
  \tag{6.5}
  \label{eq:6-5}
  \mathcal{L}_\theta=\mathcal{L}-\frac{i\theta}{2\pi}\int F,
\end{equation}

以及一个质量项：

\begin{equation}
  \tag{6.6}
  \label{eq:6-6}
  \mathcal{L}_{\theta,m}=\mathcal{L}_\theta+\frac{m}{2\pi}\bar\psi_-\psi_++\frac{\bar m}{2\pi}\bar\psi_+\psi_-,
\end{equation}

（这里 $m$ 是复的）。

\subsection{手征对称性}
% 原文 §6.2

手征对称性（chiral symmetry）是经典拉格朗日量 (6.4) 或 (6.5) 的一个 $U(1)$-对称性，由下式给出

\begin{equation}
  \tag{6.7}
  \label{eq:6-7}
  \psi_+\to\psi_+,\qquad \psi_-\to e^{i\gamma}\psi_-.
\end{equation}

这个对称性被质量项破坏。然而，即便没有质量项，由于下述拓扑理由，这个对称性也可能在量子力学上被破坏。

设时空是一张闭 Riemann 面。于是在手征对称性下，“测度” $\mu=D\psi_-D\bar\psi_-$ 变换为 $\mu\to e^{i\gamma I}\mu$，其中 $I$ 是算子 $D^A_+:S_-\otimes_{\mathbb{C}}E\to S_+\otimes_{\mathbb{C}}E$ 的指标（index）。“证明”：设 $\{\psi_-\}$ 记

% ===== 原文 p.83 =====
$S_-$ 的荷（charge）为 $1$ 的截面空间，而 $\{\bar\psi_-\}$ 记荷为 $-1$ 的对应空间。则形式上 $I=\dim(\{\psi_-\})-\dim(\{\bar\psi_-\})$；由于 $(\{\bar\psi_-\})=(\{\psi_-\})^*$，我们有 $I=\dim(\{\psi_-\})-\dim(\{\psi_+\})=\mathrm{ind}\big(D^A_+|_{S_+\otimes_{\mathbb{C}}E}\big)$。这里 $\mathrm{ind}$ 表示所述算子的指标。

% [Tutor Note: 这一步“由 $(\{\bar\psi_-\})=(\{\psi_-\})^*$ 得 $\dim(\{\bar\psi_-\})=\dim(\{\psi_+\})$”隐含使用了 $S_-\otimes S_+=\mathbb{C}$（见 p.81），即复共轭把荷 $-1$ 的 $S_-$ 截面空间同构地送到荷 $+1$ 的 $S_+$ 截面空间。若忽略这点而误以为 $\bar\psi_-$ 与 $\psi_-$ 的截面空间同维，指标将恒为零，反常就消失了。此外 $I=c_1(P)$ 是 Atiyah–Singer 指标定理在二维的显式形式（$S_\pm$ 复一维、$E$ 复一维时，$\mathrm{ind}$ 就是耦合后的度数公式）；$1/2\pi$ 这一归一化必须与 (6.5) 中 $\theta$ 项的 $1/2\pi$、以及下文 $dJ_A=F/2\pi$ 的结论取自同一套约定，任何一处改号都要三处同时改。]

已知指标 $I$ 等于 $c_1(P)=\int\frac{F_A}{2\pi}$。因此，由于丛 $P$ 在拓扑上可以非平凡，手征对称性在量子理论中被破坏。

这个效应称为反常（anomaly）：一个经典对称性在量子力学上不成立，因为测度不是不变的。与 Lecture 1 中讨论的自发对称性破缺（spontaneous symmetry breaking）的差别在于：在自发对称性破缺的情形，经典对称性在量子理论中确实存在，只是未被实现。对称性上的反常意味着该对称性根本不是量子理论的对称性。再说一遍——因为这是一点虽然初等却非常重要的地方：有反常的对称性是经典对称性，而不是量子理论的对称性；而自发破缺的对称性是量子理论的对称性，但并不保持真空不变。（与自发对称性破缺不同，反常甚至可以出现在一个时空维数下。）

下面用流（current）的语言讨论手征对称性破缺的机制。经典地，手征对称性由流

\begin{equation}
  \tag{6.8}
  \label{eq:6-8}
  J_A=\bar\psi_-\psi_-
\end{equation}

生成（这是 $(0,1)$ 型的一次形式），它当然是守恒的：$dJ_A=0$。在量子力学上，这个微分方程可以被形变：$dJ_A=\mathcal{O}$，其中 $\mathcal{O}$ 是某个取值于曲面上二次形式的算子。通过列举所有可能的算子，不难证明能出现的算子（up to a factor）只有一个：它就是曲率算子 $F$。于是 $dJ_A=\alpha F$，其中 $\alpha$ 是常数，而先前的拓扑计算表明 $\alpha=1/2\pi$。也可以通过研究 Feynman 图来建立关系 $dJ_A=F/2\pi$；这最初就是这样做的。

当存在有反常的对称性时，关联函数以非常简单的方式依赖于 $\theta$。我们将在拉格朗日量 (6.5) 的背景下说明这一陈述。设 $\mathcal{O}_i$ 是任意算子，使得在手征对称性下 $\prod_i\mathcal{O}_i\to e^{in\gamma}\prod_i\mathcal{O}_i$。则上述指标计算蕴含着：关联子 $\langle\prod\mathcal{O}_i\rangle$ 的所有非平凡贡献都来自 $c_1=n$ 的丛 $P$。（这正是“测度乘以 $\prod_i\mathcal{O}_i$ 在该对称性下不变”这一条件。）于是这个关联子具有形式 $e^{in\theta}\langle\prod\mathcal{O}_i\rangle_0$，其中 $\langle\prod\mathcal{O}_i\rangle_0$ 与 $\theta$ 无关。例如 $\langle\psi_-\psi_+^*\rangle=Ce^{i\theta}$，其中 $C$ 与 theta 无关。

若我们考虑的不是 (6.5)、而是带质量项的拉格朗日量 (6.6)，则上述论证表明

\begin{equation}
  \tag{6.9}
  \label{eq:6-9}
  \Big\langle\prod\mathcal{O}_i\Big\rangle(m,\theta)=e^{in\theta}\Big\langle\prod\mathcal{O}_i\Big\rangle_0(\tilde m),
\end{equation}

其中 $\tilde m=me^{-i\theta}$。因此，理论真正重要的参数是 $\tilde m$，我们将把它写成 $m_*e^{-i\theta_*}$，$m_*\ge0$。本讲的目标就是把这个理论作为 $\tilde m$ 的函数来研究。

% ===== 原文 p.84 =====
\subsection{带质量费米子的二维规范理论的行为}
% 原文 §6.3

现在我们要描述由拉格朗日量 (6.6) 定义的这一理论的行为，随后再为我们的论断提供依据。

首先，我们将看到：对大的 $m_*$（即 $m_*\gg e$），该理论类似于上一讲中考虑过的带质量玻色子的二维规范理论。特别地，它有质量隙（mass gap）。它在 $\theta\ne\pi$ 时具有唯一真空，而在 $\theta=\pi$ 时有两个真空。宇称（parity）或空间取向反转（$t\to t$，$x\to-x$）这一离散对称性，其作用为 $m_*\to m_*$、$\theta_*\to-\theta_*$，在 $\theta_*=0$ 或 $\pi$ 时是理论的对称性。对大的 $m_*$，若 $\theta_*=0$ 则宇称不变，若 $\theta^*=\pi$ 则宇称破缺。

接下来我们将证明：对 $\tilde m=0$，理论实际上是自由的，事实上它在一次变量代换之后成为一个自由的有质量理论。因此对小 $m_*$（即 $m_*\ll e$），理论有唯一真空与质量隙。对小 $m_*$，无论 $\theta$ 是 $0$ 还是 $\pi$，宇称都不破缺。

在 $\theta=\pi$ 以及 $m_*$ 的某个 $e$ 阶临界值处，人们相信理论会失去质量隙；事实上人们相信，在临界值 $\tilde m=-m_c$ 处理论实际上是共形的（conformal），并且同构于一个自由实费米子的理论。但这些断言的论证并不那么精确。

% [Tutor Note: 本段末尾把 $\theta_*$ 误排为 $\theta^*=\pi$（上标），与全篇的下标记号不一致，应以 $\theta_*=\pi$ 为准；$\theta_*$ 是 $\tilde m=m_*e^{-i\theta_*}$ 中的相位，其取值在 $[0,2\pi)$，与 (6.7) 中手征参数 $\gamma$ 无关，切勿混同。]

\subsection{重费米子}
% 原文 §6.4

本节我们研究 $m_*\gg e$ 的情形。此时我们可以把我们的理论视为 $e=0$ 之理论的微扰，微扰参数为 $\lambda=e/m_*$。在 $\lambda=0$ 处，我们有一个二维纯规范理论（它是自由的）与有质量费米子的自由理论的直积（为稳妥起见，我们应引入 $B=A/e$；这是合理的，因为当 $e=0$ 时只有平凡 $U(1)$-丛对路径积分有贡献）。

情形与我们上周研究过的玻色情形类似，可以用类似的论证来分析。也就是说，$e=0$ 的自由理论的小 $e$ 微扰是奇异的（singular）。形变后理论的状态空间实际上比未形变理论的状态空间更小。更确切地说，费米子禁闭（confinement）发生了：唯一允许的态（对 $\theta\ne\pi$）是总荷为 $0$ 的态（这里 $\psi_\pm$ 的荷为 $1$，$\bar\psi_\pm$ 的荷为 $-1$，且 $\theta\in[0,2\pi)$）。特别地，费米子只能成对、成四地出现，单个费米子不可能存在。

与玻色情形一样，由形变论证可知理论有质量隙（形变是奇异的这一事实并不使该论证失效，因为 Hilbert 空间不会增加，只会变小）。更确切地说，$\theta\ne\pi$ 时有一种实现，$\theta=\pi$ 时有两种（与纯规范理论一样），而任何实现都有质量隙。

\subsection{Bose–Fermi 对应}
% 原文 §6.5

在研究轻费米子的情形之前，我们先考虑 Bose–Fermi 对应（Bose–Fermi correspondence），它在研究小 $m_*$ 的情形时会有用。

考虑两个二维自由场论

% ===== 原文 p.85 =====
\begin{enumerate}
\item 由拉格朗日量
\begin{equation}
  \tag{6.10}
  \label{eq:6-10}
  \mathcal{L}_f=\frac{1}{2\pi}\int d^2x\left(\bar\psi_-\partial_+\psi_-+\bar\psi_+\partial_-\psi_+\right)
\end{equation}
定义的费米子理论。

\item 由拉格朗日量
\begin{equation}
  \tag{6.11}
  \label{eq:6-11}
  \mathcal{L}_b=\frac{R^2}{4\pi}\int d^2x\,|d\phi|^2
\end{equation}
定义的玻色子理论。
\end{enumerate}

这两个理论是共形的（经典上与量子力学上都是），且 Virasoro 中心荷为 $1$。因此我们可以猜测它们之间可能存在某种联系。而这样的联系确实存在；它被称为 Bose–Fermi 对应。

\textbf{Remark.} 中心荷很小的共形场论非常稀少。例如对 $c<1$，它们已被 Friedan、Qiu、Shenker 完全分类（Phys. Rev. Lett, vol 52, 1984, p. 1575）。答案是：不存在连续参数，理论几乎完全由 $c$ 决定，而 $c$ 只能取一个离散的值序列。对 $c=1$，这样的分类尚不可得，但已知的例子极少，且它们全都可以用自由 Bose 场来构造。

现在让我们看看玻色理论 (6.11)。如我们在 Lecture 1 中所记得的，为使这个理论有意义，$\phi$ 必须是角度值的（angle-valued），即取值于圆周 $\mathbb{R}/2\pi\mathbb{Z}$。拉格朗日量中的常数 $R$ 具有这个圆周半径的意义。

\textbf{Remark.} 我们总是把目标圆周模型化为 $\mathbb{R}/2\pi\mathbb{Z}$，但考虑其上的各种 Riemann 度量，它们由半径 $R$ 的取值来参数化。

回顾 Lecture 1：这个理论（在其唯一的实现中）的 Hilbert 空间形如 $\mathcal{H}_b=\oplus_{k\in\mathbb{Z}}(F_b\otimes F_b^*)_k$，其中 $F_b$ 是玻色 Fock 空间。理论的算子代数 $A_b$ 由 $\phi$ 的导数（但不是 $\phi$ 本身）以及 $:e^{in\phi}:$ 生成（为简便起见，我们暂时省略冒号）。算子乘积展开（operator product expansion）由 [I-Perturbative QFT] 的 Lecture 3 中的公式 (3.9) 给出，其中 $D(x-y)=-R^{-2}\ln|x-y|$。$A_b$ 在 $\mathcal{H}_b$ 中的作用为：$\phi$ 的导数不改变 $k$，而 $:e^{in\phi}:$ 把 $(F_b\otimes F_b^*)_k$ 映到 $(F_b\otimes F_b^*)_{k+n}$。

现在考虑理论 (6.10)。这个理论的 Hilbert 空间是 $\mathcal{H}_f=F_f\otimes F_f^*$，其中 $F_f$ 由真空经全纯算子 $\psi_+,\bar\psi_+$ 生成，而 $F_f^*$ 由真空经 $\psi_-,\bar\psi_-$ 生成。算子代数 $A_f$ 由 $\psi_+,\bar\psi_+,\psi_-,\bar\psi_-$ 生成，带有自由理论的标准 OPE。

更仔细地考虑算子 $:e^{in\phi}:$，$n\in\mathbb{Z}$。如我们所知（[I-Perturbative QFT] 的 Lecture 3 与 Lecture 1；该结果可由对我们刚写下的 $D(x-y)$ 公式取指数得到），这个算子具有全纯维数与反全纯维数 $n^2/4R^2$（左右维数之和为 $n^2R^2/2$，正如我们在 [I-Perturbative QFT] 的 Lecture 3 中所见）。

% [Tutor Note: 此处“左右维数之和为 $n^2R^2/2$”与前面的 $h=\bar h=n^2/4R^2$ 不能同时成立：$h+\bar h=n^2/2R^2$。由 p.86 紧接着的条件“$n^2/2R^2$ 必须是整数”（$e^{in\phi_+}$ 的单值性）可以判定，正确的和应为 $n^2/2R^2$，原文的 $n^2R^2/2$ 疑为 $R\to1/R$（T-对偶）写法残留的笔误。注意顶点算子的维数在 $R\to1/R$ 下并非不变，交换的是 $e^{in\phi}$ 与磁型/电型算子。]

经典地，算子 $e^{in\phi}$ 局部地分解为全纯因子与反全纯因子之积：$e^{in\phi}=e^{in\phi_+}e^{in\phi_-}$，其中 $\phi=\phi_++\phi_-$，且 $\partial_-\phi_+=\partial_+\phi_-=0$。（当然，这里 $\phi_+,\phi_-$ 只定义到相差一个常数。）并不显然的是，人们可以定义算子 $e^{in\phi_\pm}$，而这样的算子

% ===== 原文 p.86 =====
并不出现在代数 $A_b$ 中。但不妨设想一下 $e^{in\phi_\pm}$ 有意义。那么我们就会从 $e^{i\phi}$ 的 OPE 发现，$e^{i\phi_+}$ 具有全纯维数 $1/4R^2$ 与反全纯维数 $0$，并且

\begin{align}
  & \left\langle e^{i\phi_+(x_1)}\cdots e^{i\phi_+(x_n)}e^{-i\phi_+(y_1)}\cdots e^{-i\phi_+(y_n)}\right\rangle \tag{6.12}\label{eq:6-12}\\
  & \qquad =\frac{\prod_{1\le i<j\le n}(x_i-x_j)^{1/2R^2}\;\prod_{1\le i<j\le n}(y_i-y_j)^{1/2R^2}}{\prod_{1\le i,j\le n}(x_i-y_j)^{1/2R^2}}.\notag
\end{align}

其中 $x_i,y_j$ 被视为复数。由这个公式清楚可见，为使 $e^{i\phi_+}$ 有任何意义，我们需要 $1/2R^2\in\mathbb{Z}$，使得 (6.12) 的右端是单值的函数。类似地，为使 $e^{in\phi_+}$ 有意义，$n^2/2R^2$ 必须是一个整数。

我们考虑 $e^{in\phi_+}$ 能有意义的第一个情形，即 $R=1/\sqrt2$ 且 $n=1$。此时我们有

\begin{align}
  & \left\langle e^{i\phi_+(x_1)}\cdots e^{i\phi_+(x_n)}e^{-i\phi_+(y_1)}\cdots e^{-i\phi_+(y_n)}\right\rangle \tag{6.13}\label{eq:6-13}\\
  & \qquad =\frac{\prod_{1\le i<j\le n}(x_i-x_j)\;\prod_{1\le i<j\le n}(y_i-y_j)}{\prod_{1\le i,j\le n}(x_i-y_j)}.\notag
\end{align}

由 (6.13) 清楚可见，$e^{i\phi_+}$ 的行为像费米子；更一般地，$e^{in\phi_+}$ 在 $n$ 为奇数时行为像费米子，在 $n$ 为偶数时行为像玻色子。这使我们期望：在 $R=1/\sqrt2$ 时，我们的理论与费米子的理论有关。

让我们看看为什么确实如此。我们来计算费米子关联函数 $\langle\psi_+(x_1)\cdots\psi_+(x_n)\bar\psi_+(y_1)\cdots\bar\psi_+(y_n)\rangle$。用 Wick 公式，容易求得

\begin{equation}
  \tag{6.14}
  \label{eq:6-14}
  \left\langle\psi_+(x_1)\cdots\psi_+(x_n)\bar\psi_+(y_1)\cdots\bar\psi_+(y_n)\right\rangle=\det\left(\frac{1}{x_i-y_j}\right).
\end{equation}

(6.13)、(6.14) 右端的吻合是一个著名的组合恒等式，它来自零点与极点的比较，以及无穷远处渐近行为的比较。

\textbf{Remark.} 事实上，(6.13)、(6.14) 的显式形式与“证明它们相等”无关。有关的只是零点与极点的结构，以及无穷远处的渐近行为，而这两者在两种情形下都可由 OPE 明显看出。

设 $\widehat{A}_b$ 是由 $A_b$ 与 $e^{in\phi_\pm}$ 生成的算子代数。等式 (6.13)、(6.14) 表明存在一个同态 $\xi:A_f\to\widehat{A}_b$，由 $\xi(\psi_\pm)=e^{\pm i\phi_\pm}$、$\xi(\bar\psi_\pm)=e^{\mp i\phi_\pm}$ 定义，它保持期望值。

利用 OPE，容易求得

\begin{equation}
  \tag{6.15}
  \label{eq:6-15}
  \xi(:\bar\psi_+\psi_+ :)=\lim_{\varepsilon\to0}\xi\left(\bar\psi_+(x+\varepsilon)\psi_+(x)-\frac{1}{\varepsilon}\right)=\lim_{\varepsilon\to0}\left(e^{-i\phi_+(x+\varepsilon)}e^{i\phi_+(x)}-\frac{1}{\varepsilon}\right)=-i\partial_+\phi.
\end{equation}

% [Tutor Note: (6.15) 中被减去的 $1/\varepsilon$ 是正规序的减除项：由 $\langle\phi_+(x+\varepsilon)\phi_+(x)\rangle=-\frac{1}{2R^2}\ln\varepsilon$ 可得 $e^{-i\phi_+(x+\varepsilon)}e^{i\phi_+(x)}\sim\varepsilon^{-1/2R^2}\big(1-i\varepsilon\partial_+\phi+\cdots\big)$，在 $R=1/\sqrt2$ 时奇异因子恰为 $\varepsilon^{-1}$，减去它之后极限给出 $-i\partial_+\phi$；若 $R\ne1/\sqrt2$，减除项须改为 $\varepsilon^{-1/2R^2}$，本式不能照搬。符号上：$\xi(\psi_+)=e^{+i\phi_+}$、$\xi(\bar\psi_+)=e^{-i\phi_+}$ 这一配对方式必须与 (6.15) 中 $-i\partial_+\phi$ 的负号一致，若把 $\xi(\psi_+)$ 取成 $e^{-i\phi_+}$，则 $:\bar\psi_-\psi_-:$ 的对应符号会整体反号。]

同样地，$\xi(:\bar\psi_-\psi_- :)=i\partial_-\phi$。又 $\xi(\psi_+\bar\psi_-)=e^{i\phi}$，$\xi(\psi_-\bar\psi_+)=e^{-i\phi}$。事实上，不难看出这个同态是一个同构。

% ===== 原文 p.87 =====
在 Hilbert 空间的层面上，$\xi$ 诱导出一个双分次同构 $F_f\to F_b\otimes l_2(\mathbb{Z})$（双次数 $=$（量子标度维数，荷）），其中 $e^{i\phi_+}$ 的荷与 $\psi_+$ 的荷都等于 $1$。写出相应的特征标（character）公式，我们便得到

\begin{equation}
  \tag{6.16}
  \label{eq:6-16}
  \frac{\sum_{n\in\mathbb{Z}}q^{n^2/2}z^n}{\prod_{n\ge1}(1-q^n)}=\prod_{n\ge1}\left(1+q^{n-1/2}z\right)\left(1+q^{n-1/2}z^{-1}\right),
\end{equation}

这也就是著名的 Jacobi 三重积恒等式（Jacobi triple product identity）。（这里左端是左行玻色子的配分函数，右端是费米子的类似配分函数，所用的是圆周上的 bounding spin 结构。要 a priori 地说明这才是与玻色子比较时应采用的正确 spin 结构，需要比我们所做的更为细致的分析。）

这个对应 $\xi$ 称为 Bose–Fermi 对应。

% [Tutor Note: (6.16) 是 Jacobi 三重积恒等式的“配分函数”形式，容易与通常的 $\prod(1-q^n)(1+q^{n-1/2}z)(1+q^{n-1/2}z^{-1})=\sum_n q^{n^2/2}z^n$ 混淆：注意此处左端的分子/分母关系恰好是通常恒等式除以 $\prod(1-q^n)$，即把玻色零点能（$q^{-1/24}$ 型的 Vacuum 能量）略去不写。若计及中心荷 $c=1$ 的 Casimir 能量，两侧都还应乘 $q^{-1/24}$，但阶数比较不受影响。$z$ 的幂次度量 $U(1)$ 荷，$q$ 的幂次度量标度维数，两者共同构成双分次。]

\textbf{Remark 1.} 如果圆周的半径不是 $1/\sqrt2$ 而是 $1$，那么算子 $e^{i\phi_+}$ 没有意义，因为它的两点函数会是 $(x-y)^{-1/2}$，不是单值的。然而算子 $e^{2i\phi}$ 是有定义的，它的两点函数是 $\dfrac{1}{(x-y)^2}$。这表明 $e^{2i\phi}$ 的行为像某种对称性的流。而确实，相应的模型等价于 Kac–Moody 中心荷为 $1$ 的 $SU(2)$-WZW 模型，因此它具有 $SU(2)$ 对称性。事实上，算子 $e^{\pm2i\phi_+}$ 与 $\partial_+\phi$ 的 Fourier 分量生成左行 $\widehat{SU(2)}$，而算子 $e^{\pm2i\phi_-}$ 与 $\partial_-\phi$ 的 Fourier 分量生成右行 $\widehat{SU(2)}$。

\textbf{Remark 2.} 事实上，Bose–Fermi 对应不仅在局部（在算子乘积展开的层面上）成立，而且在整体上（在任意 Riemann 面上的路径积分层面上）也成立。Bose–Fermi 对应的一个优雅的整体版本如下。对任何闭 Riemann 面 $\Sigma$，配分函数之间有如下恒等式

\begin{equation}
  \tag{6.17}
  \label{eq:6-17}
  \int D\phi\,e^{-\frac{1}{8\pi}\int|d\phi|^2}=\sum_{\varepsilon}\int D\psi_+D\bar\psi_+D\psi_-D\bar\psi_-\,e^{-\mathcal{L}_f(\psi)},
\end{equation}

其中 $\phi$ 被理解为 $\Sigma$ 到圆周的映射，$\mathcal{L}_f$ 是由 (6.10) 给出的拉格朗日量，而 $\varepsilon$ 遍历 $\Sigma$ 上的 spin 结构（$\psi_+$ 与 $\psi_-$ 取相同的 spin 结构）。若按上面解释的方式建立算子之间的对应，则对算子的关联函数也有类似的恒等式。

\subsection{非线性理论的 Bose–Fermi 对应}
% 原文 §6.6

我们已经建立了两个自由理论——一个玻色子的理论与一个费米子的理论——之间的对应。既然我们知道如何把一侧的算子与另一侧的算子等同起来，我们就可以把“在拉格朗日量中加入给定算子以微扰玻色理论”所得的非自由理论，与“加入相应算子以微扰费米理论”所得的非自由理论相比较。这引出许多惊人的结果。我们考虑几个例子。

1. 回顾在我们的对应下 $\mathcal{L}_f\to\mathcal{L}_b(1/\sqrt2)$，其中 $\mathcal{L}_b(R)$ 由 (6.11) 给出。由于 $\xi(\bar\psi_\pm\psi_\pm)=\mp i\partial_\pm\phi$，我们得到

\begin{equation}
  \tag{6.18}
  \label{eq:6-18}
  \mathcal{L}_f+\frac{1}{2\pi}\int d^2x\left(g\bar\psi_+\psi_+\bar\psi_-\psi_-\right)\to\mathcal{L}_b\left(1/\sqrt2(1+g)^{-1/2}\right).
\end{equation}

% ===== 原文 p.88 =====
这表明，出乎我们意料，拉格朗日量为 $\mathcal{L}_f+\frac{1}{2\pi}\int d^2x(g\bar\psi_+\psi_+\bar\psi_-\psi_-)$ 的理论是自由的。特别地，它的 $\beta$ 函数为零。当理论用 Bose 变量描述时这是显然的，在 Fermi 变量下却不显然。

2. 另一方面，考虑自由有质量费米子的理论，其拉格朗日量为 $\mathcal{L}_f+\frac{1}{2\pi}\int d^2x(m\bar\psi_-\psi_++\bar m\bar\psi_+\psi_-)$。利用 $\xi(\bar\psi_\pm\psi_\mp)=-e^{\mp i\phi}$，我们得到：在 $\xi$ 之下，这个自由拉格朗日量变为

\begin{equation}
  \tag{6.19}
  \label{eq:6-19}
  \mathcal{L}_b\left(1/\sqrt2\right)-\frac{1}{2\pi}\int d^2x\left(me^{i\phi}+\bar m e^{-i\phi}\right).
\end{equation}

于是我们得到另一个惊人的事实：由 (6.19) 定义的非线性理论实际上是自由的。

\textbf{Remark.} 可能看起来 (6.19) 的第二项有错误的标度维数。并非如此，因为算子 $e^{\pm i\phi}$ 的反常维数为 $1$。更确切地说，(6.19) 并不固定一个理论，而是固定依赖于动量标度 $\mu$ 的一族理论，$\mu$ 是在算子 $e^{i\phi}$ 被重整化时引入的。这个标度出现在拉格朗日量中相应项的前面，并抵消了量纲上的不符。

利用对称性 $\phi\to\phi+\theta$，我们可以把 (6.19) 化归为 $m$ 为实的情形。此时 (6.19) 看起来是

\begin{equation}
  \tag{6.20}
  \label{eq:6-20}
  \mathcal{L}_b\left(\sqrt2\right)-\frac{m}{\pi}\int d^2x\cos\phi.
\end{equation}

% [Tutor Note: 由 (6.19) 到 (6.20) 时 $R$ 从 $1/\sqrt2$ 变成 $\sqrt2$，相当于把拉格朗日量整体乘以 $4$（因 $\mathcal{L}_b(R)\propto R^2$），此时质量项系数本应同时变为 $4m/\pi$，而原文仍写 $m/\pi$。这不影响物理（整体归一化可任意，等价于重新定义 $m$），但下面 (6.21) 用 $m'$ 与 (6.20) 作比较时，必须锁定同一个归一化约定，否则会差出因子 $4$；下文 $m'$ 的系数问题正出在同一处 $R$ 的重标度上（见下条批注）。]

于是，经典运动方程为 $\Delta\phi=-4m\sin\phi$。这个方程称为 sine-Gordon 方程，它是著名的完全可积的孤子方程。量子理论也是可积的（在 Lecture 4 中我们所了解的意义上）。当然，在量子理论等价于一个自由有质量费米子理论的那个耦合值处，它不仅是可积的，而且是自由的。

现在考虑拉格朗日量

\begin{equation}
  \tag{6.21}
  \label{eq:6-21}
  \mathcal{L}_b(R)-\frac{m}{\pi}\int d^2x\cos\phi.
\end{equation}

这个拉格朗日量与 (6.20) 成比例，其中 $m'=m/2R^2$，所以经典上两个拉格朗日量等价。然而在量子力学上并非如此，因为把拉格朗日量乘以一个常数确实会影响量子理论。事实上，理论现在本质上依赖于 $R$。若 $R=1/\sqrt2$，理论是自由的。对一般的 $R$ 它不是自由的。映射 $\xi$ 表明，对一般的 $R$，理论等价于拉格朗日量为

\begin{equation}
  \tag{6.22}
  \label{eq:6-22}
  \mathcal{L}_f+\int d^2x\left(m\bar\psi_-\psi_++m\bar\psi_+\psi_-+g\bar\psi_+\psi_+\bar\psi_-\psi_-\right),
\end{equation}

的费米理论，其中 $R=1/\sqrt2(1+g)^{-1/2}$。

% [Tutor Note: “$m'=m/2R^2$”疑为 $m'=2m/R^2$ 之误，可直接验证：$\mathcal{L}_b(R)=(R^2/2)\mathcal{L}_b(\sqrt2)$，故若 (6.21)$\,=(R^2/2)\times$(6.20)$|_{m'}$，则 $(R^2/2)(m'/\pi)=m/\pi$，即 $m'=2m/R^2$。取 $R=1/\sqrt2$（即 $g=0$）检验：此时 $m'=4m$，$(1/4)\big[\mathcal{L}_b(\sqrt2)-\frac{4m}{\pi}\int\cos\phi\big]=\mathcal{L}_b(1/\sqrt2)-\frac{m}{\pi}\int\cos\phi$，恰为 (6.21) 在 $R=1/\sqrt2$ 处的值，亦即自由费米子理论；而按原文的 $m'=m/2R^2=m$ 只会得到 $\frac{m}{4\pi}\int\cos\phi$，与“$R=1/\sqrt2$ 时理论自由”矛盾。]

% [Tutor Note: 关于 $g$ 的依赖，原文内部也不自洽：由 (6.18)、(6.22) 得 $R=(1/\sqrt2)(1+g)^{-1/2}$，即 $R^2=(1+g)^{-1}/2$，代入 (6.11) 应得 $|d\phi|^2$ 的系数 $R^2/4\pi=1/(8\pi(1+g))$；然而 (6.24) 与 (6.27) 中该系数一律写作 $(1+g)/8\pi$。二者只在 $g=0$ 时一致，故其中必有一处笔误（若以 (6.24)/(6.27) 为准，则 (6.18)、(6.22) 中的指数应为 $+1/2$）。此点直接影响“大 $g$ 对应大半径（弱耦合）”还是“大 $g$ 对应小半径”的物理判断，阅读时务必以自己重算的系数为准。]

如我们提到的，由拉格朗日量 (6.21) 描述的理论在 $R\ne1/\sqrt2$ 时不是自由的，但是可积的。它的 $S$ 矩阵可以用我们在 Lecture 4 中讨论过的方法计算。这个理论对大 $R$（即在经典

% ===== 原文 p.89 =====
极限下）的可解性，与 sine-Gordon 方程在经典层面的完全可积性有关。

现在我们想把 Bose–Fermi 对应应用于规范理论。若 $m$ 很小或为零，这将给出严格而精确的结果。考虑由 (6.6) 给出的拉格朗日量 $\mathcal{L}_{\theta,m}$。定义

\begin{equation}
  \tag{6.23}
  \label{eq:6-23}
  \mathcal{L}^f_{\theta,m}=\mathcal{L}_{\theta,m}+\frac{1}{2\pi}\int d^2x\,g\bar\psi_+\psi_+\bar\psi_-\psi_-.
\end{equation}

我们把它改写为 Bose 变量：

\begin{align}
  \mathcal{L}^b_{\theta,m} & =\int d^2x\left(\frac{1+g}{8\pi}|d\phi|^2-\frac{m}{2\pi}e^{i\phi}-\frac{\bar m}{2\pi}e^{-i\phi}\right. \tag{6.24}\label{eq:6-24}\\
  & \qquad\left.+\,A_+\left(\frac{i\partial_-\phi}{2\pi}\right)+A_-\left(\frac{-i\partial_+\phi}{2\pi}\right)+\frac{|\ast F|^2}{4e^2}-\frac{i\theta}{2\pi}F\right).\notag
\end{align}

我们对项 $\int(A_-\partial_+\phi-A_+\partial_-\phi)$ 分部积分，它等于 $\int\phi(\partial_-A_+-\partial_+A_-)$。表达式 $\partial_-A_+-\partial_+A_-$ 等于曲率 $F$，于是 (6.24) 简化为：

\begin{equation}
  \tag{6.25}
  \label{eq:6-25}
  \mathcal{L}^b_{\theta,m}=\int d^2x\left(\frac{1+g}{8\pi}|d\phi|^2+\frac{m}{2\pi}e^{i\phi}+\frac{\bar m}{2\pi}e^{-i\phi}+\frac{|\ast F|^2}{4e^2}-\frac{i(\phi+\theta)}{2\pi}F\right).
\end{equation}

% [Tutor Note: 由 (6.24) 到 (6.25) 只是把 $A_\pm$ 项分部积分化为 $\phi F$ 项，这一步同时给出 $\int(A_-\partial_+\phi-A_+\partial_-\phi)=\int\phi(\partial_-A_+-\partial_+A_-)$（符号：$\int A_-\partial_+\phi=-\int\phi\partial_+A_-$，$\int(-A_+\partial_-\phi)=+\int\phi\partial_-A_+$）以及 $\phi F$ 项与原有的 $-\frac{i\theta}{2\pi}\int F$ 合并成 $-\frac{i(\phi+\theta)}{2\pi}F$，都与原文一致。但质量项的符号在 (6.24) 中是负号、在 (6.25) 中变成正号，这不是分部积分所能产生的；它等价于同时作 $\phi\to\phi+\pi$（此时 $\phi F$ 项给出 $\theta\to\theta+\pi$ 的移动）。此处按原文照录，读者若自行推导，务必固定一种相位约定，否则 $m_*$ 与 $\theta_*$ 的相对符号会出错。]

$\phi F$ 项在整体上其实必须以 Chern–Simons 泛函的意义来定义（因为 $\phi$ 是到圆周的映射，而不是实值的），但我们不必为此小心。

由这个方程清楚可见，当 $m=0$ 时理论是自由的（因为拉格朗日量是二次的），并且没有本质的 $\theta$ 依赖（$\theta$ 可以通过把 $\phi$ 平移 $-\theta$ 而消去）。我们现在也能看到，在 $m=0$ 处（或附近）理论有质量隙。为此我们只需解经典方程。对于这样一个具有二次拉格朗日量的理论，经典方程的分析足以确定量子力学上会发生什么。我把它留给读者：请研究由上述拉格朗日量导出的经典方程，并证明 $\phi$ 满足一个有质量的波动方程。

对 $m\ne0$，通过把 $\phi$ 变为 $\phi+\theta$，我们发现 theta 角可以被吸收到 $m$ 中。这一点已在本讲开头基于指标理论与反常论证过。现在让我们看 $m\ne0$ 且很小时的动力学。由 Lecture 5 可知，对于近似为常数的外场 $\phi(x)$，

\begin{equation}
  \tag{6.26}
  \label{eq:6-26}
  \int DA\,e^{\frac{i}{2\pi}\int\phi F_A-\frac{1}{4e^2}\int|\ast F_A|^2}=C\,e^{-\frac{e^2}{2\pi}\int d^2x\min_n\left(n-\frac{\phi(x)}{2\pi}\right)^2},
\end{equation}

因此，对拉格朗日量 (6.25) 而言，$\phi$ 的有效拉格朗日量为

\begin{equation}
  \tag{6.27}
  \label{eq:6-27}
  \mathcal{L}^{\theta_*,m_*}_{\mathrm{eff}}=\int d^2x\left(\frac{1+g}{8\pi}|d\phi|^2-\frac{m_*}{\pi}\cos(\phi-\theta_*)+\frac{e^2}{2\pi}\min_n\left(n-\frac{\phi}{2\pi}\right)^2\right).
\end{equation}

% [Tutor Note: (6.26) 中的 $\min_n$ 来自二维 $U(1)$ 规范理论的瞬子（绕数）求和：$\phi$ 只在模 $2\pi$ 的意义下有定义，$n$ 标记 $\phi/2\pi$ 的整数平移（即第一 Chern 数），对 $n$ 求和后由鞍点近似给出 $\min_n$ 的形式；前面的系数 $e^2/2\pi$ 来自二维规范场传播子给出的 Coulomb 自能（$\frac{e^2}{2\pi}\ln|x-y|$）。由此也可看出 (6.28) 中 $U(\phi)$ 的两个结构：$\min_n$ 项在 $\phi\to\phi+2\pi$ 下不变，故 $\theta$ 只以 $\phi-\theta_*$ 的组合出现，$U$ 对 $\theta$ 以 $2\pi$ 为周期——这正是“$\theta$ 可被吸收进 $m$”在有效势层面的体现。当 $m_*\to0$ 时，极小点在 $\phi=\theta_*+2\pi n$ 处简并，$n$ 的选取由 $\min_n$ 项决定。]

由这个公式，我们可以分析小 $m_*$ 时的 $\theta$ 依赖。势为

\begin{equation}
  \tag{6.28}
  \label{eq:6-28}
  U(\phi)=-\frac{m_*}{\pi}\cos(\phi-\theta_*)+\frac{e^2}{2\pi}\min_n\left(n-\frac{\phi}{2\pi}\right)^2
\end{equation}

% ===== 原文 p.90 =====
对足够小的 $m_*$，它在所有 $\theta$ 下都有唯一的整体极小（模 $2\pi$），且二阶导数为正；在这个区域内，物理作为 $\theta$ 的函数是光滑的。

相反，在原规范理论的大 $m_*$ 分析中，我们发现在 $\theta=\pi$ 处有不连续性。我们无法真正把这个结果与我们从 (6.28) 学到的东西相比较，因为若 $m_*$ 不比 $e$ 小，$\phi$ 场的描述就变成强耦合的，所以只有在 $m_*$ 很小时，我们才能真正基于这个描述作出确定的陈述。

尽管如此，如果我们大胆地、在没有坚实依据的情况下假设量子理论即使对大的 $m_*$ 也可以用有效势 $U(\phi)$ 来处理，那么我们便得到一幅有吸引力的图像。对大的 $m_*$，$U(\phi)$ 的极小化在 $\theta=\pi$ 处给出不连续性，这与直接研究规范理论所得的一致。在 $m_*$ 的某个临界值处，不连续性消失。在这一点，$U(\phi)$ 的极小是退化的（二阶导数为零）。经典势的这种退化临界点，在量子理论中预期对应于一个在红外流向非平凡不动点的理论。事实上，人们相信，在两维中改变一个参数以得到这种临界点、越过它之后 $\mathbb{Z}/2\mathbb{Z}$ 对称性自发破缺，这种一般的量子行为正是“Ising 模型”的行为，而 Ising 模型等价于无质量自由实费米子的理论。这就是 §6.3 末所述猜测的依据。

\subsection*{参考文献与交叉引用}

二维规范理论在 Lecture 4 中讨论。角度值场在 $2$ 维中的自由理论，在 [II-CFT] 的 Lecture 1 与本系列的 Lecture 1 中讨论。WZW 模型在 [II-CFT] 的 Lecture 4 中讨论。

关于本讲所研究模型的更多内容，见

\begin{quote}
\textit{S. Coleman, More about the massive Schwinger model, Annals of Physics 101 (1976), 239--267.}
\end{quote}

另见（例如）Coleman 与 Weinberg 的书，它们在先前各讲的参考文献中被引用，其中有更多关于 $\theta$ 依赖与反常对称性的内容。
